% ===== tr-03 块6：第 3 章 pp.39--41 + p.42 顶部正文（至 References 标题前）=====
% 说明：定义 3.10.1 起于原文 p.38 末、主体在 p.39；为保持环境完整，整段置于本文件开头
% （p.38 的页标记留在 tr03e）。本块无编号公式。
\begin{definition}{3.10.1}
设 $B$ 是带奇点、以 $\Delta$ 为判别轨迹的热带仿射流形。设 $G$ 是带权、连通的有限图，其顶点集与边集分别记作 $G^{[0]}$ 与 $G^{[1]}$，权函数为 $w_G:G^{[1]}\to\N\setminus\{0\}$。$B$ 中的一条\emph{参数化热带曲线}是一个连续映射 $h:G\to B$，满足以下条件：
\begin{enumerate}[label=(\arabic*)]
\item 对每条边 $E\subseteq G$，$h|_{\operatorname{Int}(E)}$ 是嵌入，$h^{-1}(B_0)$ 在 $\operatorname{Int}(E)$ 中稠密，且存在截面 $u\in\Gamma(\operatorname{Int}(E),h^*(i_*\Lambda))$，它在 $h(\operatorname{Int}(E))\cap B_0$ 的每一点处与 $h(\operatorname{Int}(E))$ 相切。我们取此截面为本原的，即它不是 $h^*(i_*\Lambda)$ 的另一个截面的整数倍。
% [Tutor Note: (1) 中 $h^*(i_*\Lambda)$ 是拉回层：$i:B_0\hookrightarrow B$ 与 $h$ 复合后取整仿射函数层；$h^{-1}(B_0)$ 是集合的原像，不要与 $h^*$ 混。本原性（primitive）排除的是"整数倍"，不是"整系数线性组合"]
% [原文 p.39]
\item 对每个顶点 $v\in G^{[0]}$，设 $E_1,\dots,E_m\in G^{[1]}$ 是与 $v$ 相邻的边。设 $u_i$ 是 (1) 所许诺的 $h^*(i_*\Lambda)|_{\operatorname{Int}(E_i)}$ 的截面，取为背离 $v$ 的方向。这就定义了芽 $u_i\in h^*(i_*\Lambda)_v=(i_*\Lambda)_{h(v)}$。
\begin{enumerate}[label=(\alph*)]
\item 若 $h(v)\in B_0$，则以下\emph{平衡条件}在 $\Lambda_{h(v)}$ 中成立：
\[
\sum_{j=1}^mw_G(E_j)u_j=0.
\]
\item 若 $h(v)\notin B_0$，则以下平衡条件在 $(i_*\Lambda)_{h(v)}$ 中满足：
\[
\sum_{j=1}^mw_G(E_j)u_j=0\quad\operatorname{mod}\ (i_*\check\Lambda)^\perp_{h(v)}\cap(i_*\Lambda)_{h(v)}.
\]
\end{enumerate}
后一个群解释如下。设 $b\in B_0$ 是 $h(v)$ 附近的点；沿 $h(v)$ 与 $b$ 之间的路径作平行移动，把群 $(i_*\Lambda)_{h(v)}$ 与 $(i_*\check\Lambda)_{h(v)}$ 分别与 $\Lambda_b$、$\check\Lambda_b$ 的局部单值不变子群等同起来。于是 $(i_*\check\Lambda)^\perp_{h(v)}$ 是 $\Lambda_b$ 的子群，上述交有意义；它与 $b$ 及路径的选取无关。
\end{enumerate}
\end{definition}

于是远离 $B$ 的判别轨迹时，热带曲线的行为与 [46]、[50] 中的热带曲线十分相似，但它们可以有终止于判别轨迹上的leg。由于我们关心的是 $B$ 紧的情形，我们不希望出现跑到无穷远的leg。不过我提醒读者：这个定义是临时性的，(2)(b) 中的行为可能并非我们真正想要的。

这里我们看到一条热带椭圆曲线，实心点是判别轨迹上的点。终止于这些点的leg必须沿单值不变的方向。
% [Tutor Note: 原文此图无图号与图注，仅正文描述；图按 600 dpi 从 p.39 裁出，未重画。]
\figimg[0.62]{p3_p39.png}

现在让我们把这一点与曲线计数的问题联系起来。在退化 $\varphi:\mathcal X\to D$ 的情形，自然地考虑曲线的映射族：
\[
\begin{tikzcd}
\mathcal C \arrow[r, "f"] \arrow[d, "\pi"'] & \mathcal X \arrow[d, "\varphi"]\\
D \arrow[r, "g"'] & D
\end{tikzcd}
\]
这里 $g$ 可以是分歧覆叠，$\pi$ 是具有约化一维纤维的平坦态射。在我们关心的情形，对 $t\ne0$，$f|_{\mathcal C_t}:\mathcal C_t\to\mathcal X_t$ 应当是曲线的稳定映射。我们假设 $t\ne0$ 时 $\mathcal C_t$ 是非奇异曲线。
% [原文 p.40]
在对数背景下，于是自然把由除子 $\mathcal C_0\subseteq\mathcal C$ 诱导的对数结构放到 $\mathcal C$ 上，从而得到对数态射的交换图
\[
\begin{tikzcd}
\mathcal C^\dagger \arrow[r, "f"] \arrow[d, "\pi"'] & \mathcal X^\dagger \arrow[d, "\varphi"]\\
D^\dagger \arrow[r, "g"'] & D^\dagger
\end{tikzcd}
\]
限制到中心纤维，我们得到交换图
\[
\begin{tikzcd}
\mathcal C_0^\dagger \arrow[r, "f"] \arrow[d, "\pi"'] & \mathcal X_0^\dagger \arrow[d, "\varphi"]\\
0^\dagger \arrow[r, "g"'] & 0^\dagger
\end{tikzcd}
\]
这提示我们应当建立稳定对数映射与对数 Gromov--Witten 不变量的理论。这一理论应当推广 Li--Ruan [42] 与 Jun Li [41] 所发展的理论。我在这里对此谈得很少，因为它很快会变得相当技术化。Siebert 在这方面有正在进行的工作。这个观点已经在 [48] 中用于环面簇中的曲线计数，因此在那里可以找到这一途径的更多线索。作为替代，我想勾画一下这样的图表如何产生一条热带曲线。
% [Tutor Note: 三个交换图的竖直箭头都是 $\pi$（左）与 $\varphi$（右），水平箭头皆为 $f$（上）与 $g$（下）；底行左右两端都是 $D$（中心纤维情形为 $0^\dagger$），不要误把左下角写成 $\mathcal C$ 之类]

为此，考虑 $\pi$ 为正规交叉、且诱导映射 $\mathcal C_0^\dagger\to\mathcal X_0^\dagger$ 在 $0^\dagger$ 上没有无穷小对数自同构的情形。（这是稳定映射这一概念的对数版本。）设 $(B,\mathcal P)$ 是对数 Calabi--Yau 空间 $\mathcal X_0^\dagger$ 的对偶相交复形。我们可以定义 $f:\mathcal C_0^\dagger\to\mathcal X^\dagger$ 的\emph{对偶相交图}，它将是 $B$ 上的一条参数化热带曲线。这里我只做 $f$ 的像与定义 3.7.1 的 (4) 中集合 $Z\subseteq\mathcal X$ 不相交的情形；否则还有一些技术细节需要处理。

首先构造 $G$。设 $C_1,\dots,C_m$ 是 $\mathcal C_0$ 的不可约分支。为便于描述这一构造，假设这些分支都是正规的。置 $G^{[0]}=\{v_1,\dots,v_m\}$。另一方面，只要 $C_i\cap C_j\ne\varnothing$，$G^{[1]}$ 就含有一条连接 $v_i$ 与 $v_j$ 的边 $\overline{v_iv_j}$。

为定义 $h:G\to B$，我们先描述每个顶点的像。设 $X_{\sigma_i}$ 是 $\mathcal X_0$ 中包含 $f(C_i)$ 的最小分层，其中 $\sigma_i\in\mathcal P$。设 $\eta_i$ 是 $C_i$ 的一般点，$\xi_i=f(\eta_i)$。由于 $f$ 是对数态射，我们有诱导映射 $f^\#:\mathcal M_{\mathcal X_0,\xi_i}\to\mathcal M_{\mathcal C_0,\eta_i}$。这在幽灵层的茎上诱导出一个交换图
\[
\begin{tikzcd}
\overline{\mathcal M}_{\mathcal C_0,\eta_i} & \overline{\mathcal M}_{\mathcal X_0,\xi_i} \arrow[l, "\bar f^\#"']\\
\overline{\mathcal M}_0 \arrow[u, "\bar\pi^\#"] & \overline{\mathcal M}_0 \arrow[l, "\bar g^\#"] \arrow[u, "\bar\varphi^\#"']
\end{tikzcd}
\]
现在 $\overline{\mathcal M}_0=\N$（见例 3.8.2 的 (4)），且由于 $\pi$ 是正规交叉，$\overline{\mathcal M}_{\mathcal C_0,\eta_i}=\N$。另一方面，$\bar\pi^\#$ 是恒等映射；而若 $g$ 是次数为 $d$ 的分歧覆叠，则 $\bar g^\#$ 是乘以 $d$。由习题 3.8.3，
\[
\overline{\mathcal M}_{\mathcal X_0,\xi_i}=\operatorname{Hom}_{monoid}(C(\sigma_i)\cap(M\oplus\Z),\N).
\]
% [Tutor Note: 上图的水平箭头都朝左（方向由 $\bar f^\#$、$\bar g^\#$ 的 $Y\leftarrow X$ 记号决定），竖直箭头都朝上；tikz 的 "leftarrow" 需写成 \arrow[l]，别写成 \arrow[r]]

% [原文 p.41]
但
\[
\operatorname{Hom}_{monoid}(\operatorname{Hom}_{monoid}(C(\sigma_i)\cap(M\oplus\Z),\N),\N)=C(\sigma_i)\cap(M\oplus\Z),
\]
所以 $\bar f^\#$ 由一个元素 $(m,r)\in C(\sigma_i)\cap(M\oplus\Z)$ 决定。现在
\[
\bar f^\#(\bar\varphi^\#(1))=\bar f^\#(0,1)=\langle(m,r),(0,1)\rangle=r,
\]
而
\[
\bar\pi^\#(\bar g^\#(1))=\bar\pi^\#(d)=d.
\]
于是 $r=d$，且 $m/d\in\sigma_i$。我们定义 $h(v_i)=m/d$。这是 $\sigma_i$ 中一点，因而含于 $B$ 中。

若 $C_i\cap C_j\ne\varnothing$，则存在包含 $C_i\cap C_j$ 的最小分层 $X_{\sigma_{i,j}}$。当然 $X_{\sigma_{i,j}}\subseteq X_{\sigma_i}\cap X_{\sigma_j}$。特别地，$\sigma_{i,j}$ 包含 $\sigma_i$ 与 $\sigma_j$。我们取 $h(\overline{v_iv_j})$ 为 $\sigma_{i,j}$ 中连接 $h(v_i)$ 与 $h(v_j)$ 的直线段。此外，若把 $\sigma_{i,j}\subseteq M_\R$ 嵌入为格多面体，设 $m_{ij}$ 是平行于 $m_i-m_j$ 的本原格元素，并用方程
\[
w_G(\overline{v_iv_j})m_{ij}=\#(C_i\cap C_j)(m_i-m_j)
\]
来定义 $w_G(\overline{v_iv_j})$。

\begin{proposition}{3.10.2}
$h$ 是一条参数化热带曲线。
\end{proposition}

这里我们不给出证明。$X_{\sigma_i}$ 总是 $\mathcal X_0$ 的不可约分支的情形已基本上包含在 [48] 中。作为替代，我们来做另一个极端情形，它能体现对数几何的一些有趣特征。

\begin{example}{3.10.3}
设 $\mathcal C_0$ 的一个分支 $C_1$ 以及所有与 $C_1$ 相交的分支 $C_2,\dots,C_t$ 都被 $f$ 映到 $\mathcal X_0$ 的某个零维分层 $X_\sigma$。不妨设
\[
\mathcal X_0=V(\sigma)=\operatorname{Spec}\C[C(\sigma)^\vee\cap(N\oplus\Z)]/(z^{(0,1)}),
\]
如第 8 节所定义。于是 $h$ 把 $v_1,\dots,v_t$ 映到点 $m_1/d,\dots,m_t/d\in\sigma$。让我们弄清平衡条件为何在 $m_1/d$ 处成立。设 $U\subseteq\mathcal C_0$ 是 $C_1$ 的一个开邻域，它只与 $C_1,\dots,C_t$ 相交，于是作为通常态射（而非对数态射），$h$ 在 $U$ 上是常值的。把 $\mathcal C_0$ 上的对数结构限制到 $U$。我们有正合列
\[
1\to\mathcal O_U^\times\to\mathcal M_U^{\mathrm{gp}}\xrightarrow{p}\overline{\mathcal M}_U^{\mathrm{gp}}\to0.
\]
取整体截面，得到
\[
1\to\Gamma(U,\mathcal O_U^\times)\to\Gamma(U,\mathcal M_U^{\mathrm{gp}})\xrightarrow{p}\Gamma(U,\overline{\mathcal M}_U^{\mathrm{gp}})=\Z^t\xrightarrow{q}\operatorname{Pic}U.
\]
截面 $s\in\Gamma(U,\overline{\mathcal M}_U^{\mathrm{gp}})$ 定义一个 $\mathcal O_U^\times$-挠子 $p^{-1}(s)$，它在 $U$ 的 Picard 群中的类为 $q(s)$。对数几何中的一个简单练习是证明：若 $s$ 是 $\Z^t$ 的第 $i$ 个标准基向量，则 $q(s)=\mathcal O_{\mathcal C}(-C_i)|_U$。注意对 $i=2,\dots,t$ 有 $\deg\mathcal O_{\mathcal C}(-C_i)|_{C_1}=-\#C_1\cap C_i$，而 $\deg\mathcal O_{\mathcal C}(-C_1)|_{C_1}=\sum_{i=2}^t\#C_1\cap C_i$，这是因为 $\mathcal C$ 中 $C_1.C_0=0$。
% [Tutor Note: 这里的 $\#$ 表示集合的基数（即交点数）；两处度数分别是 $\deg\mathcal O_{\mathcal C}(-C_i)|_{C_1}$ 与 $\deg\mathcal O_{\mathcal C}(-C_1)|_{C_1}$，限制的底一个在 $C_1$、一个在 $C_1$ 而线丛指标分别为 $i$ 与 $1$，抄写时勿混]

现在注意，$f^\#$ 作用在幺半群层上，诱导出交换图
{\small
\[
\begin{tikzcd}
\mathcal O^\times_{V(\sigma),x}\oplus\overline{\mathcal M}^{\mathrm{gp}}_{V(\sigma),x} \arrow[r, "\cong"] & \mathcal M^{\mathrm{gp}}_{V(\sigma),x} \arrow[r, "f^\#"] \arrow[d] & \Gamma(U,\mathcal M_U^{\mathrm{gp}}) \arrow[d, "p"]\\
N\oplus\Z \arrow[r, "\cong"'] & \overline{\mathcal M}^{\mathrm{gp}}_{V(\sigma),x} \arrow[r, "\bar f^\#"'] & \Gamma(U,\overline{\mathcal M}_U^{\mathrm{gp}}) \arrow[r, "\cong"'] & \Z^t
\end{tikzcd}
\]
}
% [Tutor Note: 该图上下两行的箭头标注位置不同：上行 $\cong$ 与 $f^\#$ 标在箭头上方，下行 $\cong$、$\bar f^\#$、$\cong$ 标在箭头下方（原文如此），$p$ 标在竖直箭头右侧。两条竖直箭头分别接在第二列（$M^{\mathrm{gp}}\to\overline{M}^{\mathrm{gp}}$，无标签）与第三列（$\Gamma(U,M_U^{\mathrm{gp}})\to\Gamma(U,\overline{M}_U^{\mathrm{gp}})$，标 $p$），第一列下方没有竖直箭头；抄写时勿把左边的竖直箭头接到第一列]
\end{example}

% [原文 p.42]
映射 $\bar f^\#$ 按构造由
\[
(n,r)\in N\oplus\Z\mapsto(\langle(n,r),(m_i,d)\rangle)_{i=1,\dots,t}
\]
给出。另一方面，为使 $\bar f^\#$ 能提升为 $f^\#$，$\mathcal O_U^\times$-挠子 $p^{-1}(\bar f^\#(n,r))$ 必须对每个 $(n,r)\in N\oplus\Z$ 都有截面，即在 Picard 群中是平凡的。这蕴涵
\[
\deg\bigotimes_{i=1}^t(\mathcal O_{\mathcal C}(-C_i)|_{C_1})^{\otimes\langle(n,r),(m_i,d)\rangle}=0,
\]
亦即
\[
\sum_{i=2}^t(\#C_1\cap C_i)(\langle(n,r),(m_i,d)-(m_1,d)\rangle)=0
\]
对一切 $(n,r)\in N\oplus\Z$ 成立。但这等价于
\[
\sum_{i=2}^t(\#C_1\cap C_i)(m_i-m_1)=0,
\]
这正是平衡条件。\proofend

遵循镜像对称的逻辑，这暗示锥侧的热带曲线应当与周期有关。直到最近，人们才开始对这一点有所理解；遗憾的是，我没有篇幅与时间来详述。我只说：在 [23] 中，Siebert 与我给出了问题 3.8.7 之 (2) 的一个解答，前提是对 $X_0(B,\mathcal P,s)^\dagger$ 作某些假设，而这些假设由 $B$ 的单性蕴涵。在这个解答中，我们逐阶构造对数 Calabi--Yau 空间的显式形变。形式上，我们的构造与 Kontsevich--Soibelman [39] 从仿射流形构造非阿基米德 K3 曲面的做法有几分相似，并且用到了 [39] 的一个关键引理。然而，Kontsevich 与 Soibelman 是在我们所谓的扇侧工作，而我们是在锥侧工作。考虑到我们关于"构造对数 Calabi--Yau 空间"策略的全部讨论都是在扇侧进行的，这或许令人惊讶。但是，若认真对待镜像哲学、并希望在光滑化的描述中看到热带曲线出现，我们就必须在锥侧工作。结果发现这极其自然。事实上，所有热带有理曲线都在我们的构造中起作用。最终，所有周期都可以用我们构造所涉及的数据计算出来；特别地，周期的计算与 $B$ 上热带有理曲线的存在之间有明确的关系。一旦这一点被完全理解，就能最终牢固地理解镜像对称的几何解释。
